% =====================================================================
%  Thème 3 — EXERCICES — Niveau MOYEN
% =====================================================================
\documentclass[11pt,a4paper]{article}
\usepackage[utf8]{inputenc}
\usepackage[T1]{fontenc}
\usepackage{lmodern}
\usepackage[french]{babel}
\usepackage{amsmath,amssymb,mathtools}
\usepackage{icomma}
\usepackage{xcolor}
\usepackage[most]{tcolorbox}
\usepackage{enumitem}
\usepackage{array}
\usepackage{booktabs}
\usepackage{fancyhdr}
\usepackage[hidelinks]{hyperref}
\usepackage[a4paper,margin=2.2cm,headheight=26pt,headsep=10pt]{geometry}

\definecolor{primary}{RGB}{223,87,99}
\definecolor{primarydark}{RGB}{150,42,55}

\tcbset{boxbase/.style={enhanced,breakable,boxrule=0.9pt,arc=2.5pt,
   left=3mm,right=3mm,top=2.3mm,bottom=2.3mm,
   fonttitle=\bfseries,coltitle=white,toptitle=0.6mm,bottomtitle=0.6mm}}
\newtcolorbox[auto counter]{exercice}[1][]{boxbase,colback=primary!4,colframe=primary,
  title={\textsc{Exercice}~\thetcbcounter\ifx\relax#1\relax\relax\else\textmd{ --- \itshape #1}\fi}}

\pagestyle{fancy}
\fancyhf{}
\fancyhead[L]{\small\textcolor{primarydark}{\textbf{Fiche 6} — Les limites}}
\fancyhead[R]{\small\textcolor{primarydark}{\textsc{Thème 3 — Moyen}}}
\fancyfoot[C]{\small\thepage}
\renewcommand{\headrulewidth}{0.8pt}
\renewcommand{\headrule}{\hbox to\headwidth{\color{primary}\leaders\hrule height \headrulewidth\hfill}}

\newcommand{\R}{\mathbb{R}}

\begin{document}

\begin{center}
{\Large\bfseries\color{primarydark}Thème 3 — Limites à l'infini : le terme dominant}\\[2pt]
{\large\color{primary}Exercices — Niveau Moyen}
\end{center}
\vspace{4pt}
{\small\itshape On mélange les trois cas et les deux côtés $\pm\infty$, puis on relie le résultat aux
asymptotes horizontales.}
\vspace{6pt}

\begin{exercice}[le même polynôme des deux côtés]
Soit $P(x)=-2x^3+4x^2-1$.
\begin{enumerate}[label=\alph*),leftmargin=2em,itemsep=3pt]
  \item Calculer $\displaystyle\lim_{x\to+\infty} P(x)$.
  \item Calculer $\displaystyle\lim_{x\to-\infty} P(x)$.
  \item Les deux résultats sont-ils opposés ? Pourquoi (parité du degré) ?
\end{enumerate}
\end{exercice}

\begin{exercice}[rationnelles : les trois cas]
Déterminer chaque limite en précisant le cas ($n>m$, $n=m$ ou $n<m$).
\begin{enumerate}[label=\alph*),leftmargin=2em,itemsep=3pt]
  \item $\displaystyle\lim_{x\to+\infty}\frac{6x^2-x+2}{2x^2+5}$
  \item $\displaystyle\lim_{x\to-\infty}\frac{x^2+3}{x^3-x}$
  \item $\displaystyle\lim_{x\to-\infty}\frac{4x^3-x}{2x^2+1}$
\end{enumerate}
\end{exercice}

\begin{exercice}[asymptote horizontale]
Soit $f(x)=\dfrac{5x^2-3}{2x^2+x+1}$.
\begin{enumerate}[label=\alph*),leftmargin=2em,itemsep=3pt]
  \item Calculer $\displaystyle\lim_{x\to+\infty} f(x)$ et $\displaystyle\lim_{x\to-\infty} f(x)$.
  \item En déduire l'équation de l'asymptote horizontale de la courbe de $f$.
\end{enumerate}
\end{exercice}

\begin{exercice}[justification par mise en facteur]
On veut prouver $\displaystyle\lim_{x\to+\infty}\frac{3x^2+x}{x^2-4}=3$ sans la règle rapide.
\begin{enumerate}[label=\alph*),leftmargin=2em,itemsep=3pt]
  \item Mettre $x^2$ en facteur au numérateur et au dénominateur.
  \item Simplifier, puis passer à la limite (chaque $\frac{\cdot}{x^k}\to 0$) et conclure.
\end{enumerate}
\end{exercice}

\begin{exercice}[attention au signe du côté $-\infty$]
Calculer, en soignant les signes.
\begin{enumerate}[label=\alph*),leftmargin=2em,itemsep=3pt]
  \item $\displaystyle\lim_{x\to-\infty}\frac{x^4-1}{-2x^2}$
  \item $\displaystyle\lim_{x\to-\infty}\frac{3x^5}{x^4+1}$
\end{enumerate}
\end{exercice}

\end{document}
